% from ega4-15-terminal-fr.tex
\oldpage[IV-3]{251}
\phantomsection
\addcontentsline{toc}{section}{Index terminologique}
\section*{INDEX TERMINOLOGIQUE}

\begingroup
\footnotesize
\setlength{\parindent}{0pt}
\setlength{\parskip}{1.1pt}
\newcommand{\EGAIVthreeTerm}[2]{#1 : #2.\par}

\EGAIVthreeTerm{Critère de platitude par fibres}{11.3.10}
\EGAIVthreeTerm{Espace algébrique sur $k$, espace $k$-algébrique, espace
préalgébrique sur $k$, espace $k$-préalgébrique}{10.10.2}
\EGAIVthreeTerm{Espace de Jacobson}{10.3.1}
\EGAIVthreeTerm{Espace $k$-préalgébrique réduit}{10.10.4}
\EGAIVthreeTerm{Famille de morphismes schématiquement dominante}{11.10.2}
\EGAIVthreeTerm{Famille de morphismes universellement schématiquement dominante}{11.10.7}
\EGAIVthreeTerm{Famille de sous-préschémas schématiquement dense}{11.10.2}
\EGAIVthreeTerm{Famille séparante d'homomorphismes de faisceaux de modules}{11.9.1}
\EGAIVthreeTerm{Famille séparante d'homomorphismes de modules}{11.9.4}
\EGAIVthreeTerm{Famille universellement séparante d'homomorphismes de faisceaux de modules}{11.9.14}
\EGAIVthreeTerm{Géométriquement isomorphes ($k$-préschémas)}{9.1.4}
\EGAIVthreeTerm{Homomorphisme de Modules universellement injectif}{11.9.14 et 11.9.18}
\EGAIVthreeTerm{Lemme de Chow pour les morphismes de présentation finie}{8.10.5.1}
\EGAIVthreeTerm{« Main theorem » de Zariski pour les morphismes de présentation finie}{8.12.6}
\EGAIVthreeTerm{Morphisme équidimensionnel en un point, morphisme équidimensionnel}{13.2.2 et 13.3.2}
\EGAIVthreeTerm{Morphisme propre en un point}{15.7.1}
\EGAIVthreeTerm{Morphisme pseudo-fini}{8.12.3}
\EGAIVthreeTerm{Morphisme submersif, morphisme universellement submersif}{15.7.8}
\EGAIVthreeTerm{Morphisme universellement ouvert en un point}{14.3.3}
\EGAIVthreeTerm{Ouvert ultraaffine}{10.9.5}
\EGAIVthreeTerm{Partie finie saturée d'un $k$-préschéma algébrique}{9.8.9}
\EGAIVthreeTerm{Partie propre en un point (d'un $S$-préschéma)}{15.7.1}
\EGAIVthreeTerm{Partie quasi-constructible, partie localement quasi-constructible d'un espace topologique}{10.1.1}
\EGAIVthreeTerm{Partie très dense d'un espace topologique}{10.1.3}
\EGAIVthreeTerm{Polynôme géométriquement irréductible}{9.7.4}
\EGAIVthreeTerm{Préschéma de Jacobson}{10.4.1}
\EGAIVthreeTerm{Préschéma équidimensionnel sur un autre en un point, équidimensionnel sur un autre}{13.2.2 et 13.3.2}
\EGAIVthreeTerm{Préschéma essentiellement affine sur un autre}{8.13.4}
\EGAIVthreeTerm{Principe de l'extension finie}{9.1.1}
\EGAIVthreeTerm{Profondeur rectifiée}{10.8.1}
\EGAIVthreeTerm{Pro-objet, morphisme de pro-objets}{8.13.3}
\EGAIVthreeTerm{Pro-objet essentiellement affine}{8.13.4, 8.14.1}
\EGAIVthreeTerm{Propriété constructible, propriété ind-constructible}{9.2.1 et 9.2.2, (i) et (ii)}
\EGAIVthreeTerm{Quasi-homéomorphisme d'espaces topologiques}{10.2.2}
\EGAIVthreeTerm{Quasi-isomorphisme d'espaces annelés}{10.2.8}
\EGAIVthreeTerm{Saturée d'une partie finie d'un $k$-préschéma algébrique}{9.8.9}
\EGAIVthreeTerm{Sous-ensemble d'un $S$-préschéma submersif, universellement submersif sur son image}{15.7.8}
\EGAIVthreeTerm{Squelette primaire d'un Module, squelette virtuel}{9.8.9}
\EGAIVthreeTerm{Spectre maximal d'un anneau de Jacobson}{10.9.3}
\EGAIVthreeTerm{Type primaire d'un Module, type primaire géométrique d'un Module}{9.8.9}
\EGAIVthreeTerm{Ultrapréschéma, morphisme d'ultrapréschémas}{10.9.5}
\EGAIVthreeTerm{$R$-ultrapréschéma}{10.10.2}

\let\EGAIVthreeTerm\relax
\endgroup

\clearpage
% La page imprimée 252 est blanche.
\thispagestyle{empty}
\null

\clearpage
